\section{The saturating-removal model}\label{sec:model}

\paragraph{Dynamics.} Each individual carries a scalar damage level $X(t)\ge 0$ that evolves as
\begin{equation}\label{eq:sr}
  \dd X = \Bigl(\eta\,Y(t) - \beta\,\frac{X}{X+\kappa}\Bigr)\dd t + \sqrt{2\eps}\,\dd W_t,
  \qquad \frac{\dd Y}{\dd t} = A, \qquad X(0)=0,
\end{equation}
and dies at the first time $X$ reaches the threshold $\xc$ \citep{karin2019senescent,raz2026damage}. Damage is produced at
a rate that rises linearly with age ($\eta$, the production slope), removed by a saturating process with maximal rate
$\beta$ and half-saturation $\kappa$, and perturbed by white noise of strength $\eps$. In an untreated individual
$Y(t)=t$ ($A=1$). An intervention may change $A$, the \emph{DPU (damage production unit) accumulation rate}: it then
speeds or slows production only, while removal, noise and time itself are unchanged. An optional age-independent
external hazard $h$ (a Makeham term \citep{makeham1860law}) adds deaths that do not pass through $X$.

\paragraph{Units and the threshold convention.} A survival curve is unchanged under $X\to X/\xc$ with
$(\eta,\beta,\kappa,\eps)\to(\eta/\xc,\beta/\xc,\kappa/\xc,\eps/\xc^2)$, so $\xc$ is not identifiable from survival
data and is fixed at $\xc=1$; damage is measured in units of the threshold. The fitted parameters are
$\theta=(\eta,\beta,\eps,\kappa)$, with $\eta$ in units of $\text{time}^{-2}$ and $\beta,\eps$ in $\text{time}^{-1}$.
Two derived quantities recur below: the mean removal over damage levels below threshold,
\begin{equation}\label{eq:rbar}
  \beta\,\rbar(\kappa), \qquad \rbar(\kappa) = \int_0^1 \frac{x}{x+\kappa}\,\dd x = 1-\kappa\ln\!\Bigl(1+\frac1\kappa\Bigr),
\end{equation}
and the ballistic score $B=\beta\,\rbar(\kappa)/(\eta L)$, the mean removal capacity relative to production at the
lifespan scale $L$ \citep{raz2026damage}.

\paragraph{Time rescaling is not a change of $A$.} Speeding up an individual's whole life by a factor $a$,
$X_a(t):=X(a t)$, gives
$\dd X_a=(a^2\eta\,t - a\beta\,X_a/(X_a+\kappa))\,\dd t+\sqrt{2a\eps}\,\dd W_t$, i.e.
$(\eta,\beta,\eps,\kappa)\to(a^2\eta,a\beta,a\eps,\kappa)$: every rate changes. We call this \emph{time rescaling}.
It is a different hypothesis from a change of $A$, which multiplies only production. A change of $A$ that starts at
birth is the same as a change of $\eta$ (production $\eta A t$); the two differ only when the change starts at an age
$t_0>0$, where a change of $A$ keeps production continuous and changes its slope, whereas a change of $\eta$ makes
production jump.

\paragraph{From data to claims.} Figure~\ref{fig:pipeline} summarises how the model is fitted to a cohort and what
is reported (Sections~\ref{sec:likelihood}--\ref{sec:effects}).

% Fig. 1: the pipeline (TikZ; no data).
\begin{figure}[t]
\centering
\begin{tikzpicture}[
  font=\small,
  box/.style={draw, rounded corners=2pt, align=left, inner sep=4pt, text width=#1},
  box/.default=4.4cm,
  hd/.style={font=\small\bfseries},
  ar/.style={-{Stealth[length=2mm]}, thick},
  node distance=5mm and 5mm]
\node[box] (data) {\textbf{Data}\\[1pt]
  ages at death or censoring\\ recording resolution $\delta$\\ optional: entry age, window};
\node[box, right=of data] (lik) {\textbf{Exact likelihood} (\S\ref{sec:likelihood})\\[1pt]
  Fokker--Planck equation for $X$\\ absorbed flux gives $S(t)$, $f(t)$\\ $\ell(\theta)$ for recorded, censored\\
  and truncated data};
\node[box, right=of lik] (cal) {\textbf{Calibrator} (\S\ref{sec:estimation})\\[1pt]
  seed, multistart L-BFGS-B\\ basin re-fit, profile likelihood\\ variants J, H, W, kfix, W0};
\node[box, below=of cal] (claims) {\textbf{Claims} (\S\ref{sec:identifiability})\\[1pt]
  a value, a bound or none for\\ $\eta$, $\eps$, $\sqrt\eta/\eps$, $\beta\rbar(\kappa)/\eta$\\ $\kappa$ is never reported};
\node[box, below=of lik] (eff) {\textbf{Group comparison} (\S\ref{sec:effects})\\[1pt]
  $M_0$, singles, full model\\ naming rule (seven conditions)\\ gives a parameter, a confusion\\ set or a block};
\node[box, below=of data] (val) {\textbf{Validation} (\S\ref{sec:validation})\\[1pt]
  numerical references\\ \C{v6-n-cohorts} synthetic cohorts\\ recovery simulations};
\draw[ar] (data) -- (lik);
\draw[ar] (lik) -- (cal);
\draw[ar] (cal) -- (claims);
\draw[ar] (lik) -- (eff);
\draw[ar] (claims) -- (eff);
\draw[ar, dashed] (val) -- (eff);
\draw[ar, dashed] (val.north east) -- (lik.south west);
\end{tikzpicture}
\caption{\textbf{From survival data to claims about SR parameters.} A cohort's ages at death (with censoring, recording
resolution and an optional entry age or fit window) enter an exact likelihood computed from the Fokker--Planck equation
of the model. The calibrator maximises it and reports, for each identified quantity, a value, a bound or no claim. Two
groups are compared by fitting effect models on the same likelihood. Each part is validated against numerical
references, a library of synthetic cohorts with known parameters, and recovery simulations (dashed).}
\label{fig:pipeline}
\end{figure}


\paragraph{Interventions and group differences.} A treatment or a difference between groups is represented as a
multiplier $m$ on one parameter $p\in\{\eta,\beta,\eps,\kappa,\xc,A,h\}$, or on time (time rescaling), acting on an age
window $[t_0,t_1)$. Before $t_0$ the treated group shares the reference parameters exactly. $t_0=0$ compares two
cohorts over the whole life; $t_0>0$ models a late-life intervention. A threshold change at $t_0>0$ is implemented by
moving the absorbing boundary, not by the rescaling above, because the rescaling is exact only from birth
(Section~\ref{sec:val-effects}).
